% =====================================================================
% Encoder fine-tuning study. Drop into paper.tex with
%   % =====================================================================
% Encoder fine-tuning study. Drop into paper.tex with
%   % =====================================================================
% Encoder fine-tuning study. Drop into paper.tex with
%   % =====================================================================
% Encoder fine-tuning study. Drop into paper.tex with
%   \input{sec_finetuning}
% placed AFTER \section{Internal Comparison: Cross-Validation}
% (\label{sec:internal}) and BEFORE \section{Discussion and Limitations}.
% Iteration 1 is complete on the development grid; the official-test and
% cross-validation blocks marked %TODO are filled once cluster jobs 3/4 land.
% =====================================================================
\section{Does Fine-Tuning the Encoder Help?}
\label{sec:finetune}

Every result so far keeps the sentence encoder \emph{frozen}: the
single-lever study (\Cref{tab:levers}) showed the remaining headroom is in the
\emph{ranking}, not the decision rule, therefore to lift the AUC ceiling a possibility is
to fine-tune the encoder.
Fine-tuning an $n\!=\!107$ corpus is the regime where over-fitting is most
likely, so the study is run as a pre-registered, staged search that selects on
development data and touches the official test split only once, for a small set
of finalists.

\subsection{Fine-tuning menu and staged protocol}
\label{sec:finetune-protocol}
We compare five adaptation regimes against the frozen baseline, ordered by
trainable-parameter budget: (i) \textbf{frozen} (control, embeddings only);
(ii) \textbf{BitFit} (bias terms only); (iii) \textbf{LoRA} low-rank adapters on
the attention (optionally FFN) projections, rank $r\!\in\!\{8,16,32\}$;
(iv) \textbf{last-$k$} unfreezing of the top $k\!\in\!\{2,4\}$ transformer
blocks with layer-wise LR decay; and (v) \textbf{full} fine-tuning. We also test
\textbf{domain-adaptive pre-training} (DAPT): a masked-LM pass over the training
interviews before the supervised stage. The encoder is \texttt{bge-large-en-v1.5}
throughout; one arm swaps in \texttt{mxbai-embed-large-v1} to check
encoder sensitivity. The temporal GRU pooling head
(\Cref{sec:method}) is unchanged.

Selection is staged to ring-fence the test split:
\begin{itemize}
  \item \textbf{Stage 1 --- development grid.} 20 configurations, 5 seeds each,
        \emph{official protocol but no test evaluation}. Configurations are ranked on the 5-seed
        \emph{per-seed mean} development ROC-AUC --- the same quantity reported
        at test time.
  \item \textbf{Stage 2 --- finalists.} Every configuration whose per-seed mean
        development AUC is within one seed-noise band ($\pm0.015$) of the best,
        capped at four, plus the frozen control, advances. Each gets a
        leakage-free decision threshold from a nested out-of-fold (OOF) probe and
        is then evaluated \emph{once} on the official test split over 30 seeds.
  \item \textbf{Stage 3 --- cross-validation.} The single winner and the frozen
        control are re-evaluated under $K$-fold (with a repeat) and
        Monte-Carlo cross-validation on the training pool, mirroring
        \Cref{sec:internal}.
\end{itemize}
All encoder adaptation happens \emph{inside} each fold; no test or
held-out interview is ever seen during fine-tuning or thresholding.

\subsection{Development-grid selection}
\label{sec:finetune-iter1}

A practical caveat governs how the grid is read. The development split is
$n\!=\!35$ interviews (12 positive); at this size the \emph{thresholded} metrics
are quantised --- many unrelated configurations land on the identical confusion
matrix --- so they carry little ranking signal. We therefore rank on the
per-seed mean ROC-AUC, the only continuous discriminator, and treat AUC gaps
below $\approx\!0.01$ as noise. \Cref{tab:ft-grid} reports the representative
configurations.

These dev figures are produced by the fine-tuning harness --- which trains the
MIL head with a shorter, more strongly regularised schedule than the
design-ablation pipeline of \Cref{tab:abldesign-rep} (and embeds under mixed
precision) --- so they are internally consistent but \emph{not} numerically
comparable to that table's frozen dev AUC. All fine-tuning effects are therefore
measured against the in-harness frozen control, not against the design-ablation
numbers; the headline frozen reference remains the once-only official-test bar.

\begin{table}[t]
\caption{Fine-tuning grid, development split, ranked by the
selection metric. \textbf{AUC} is the per-seed mean$\pm$std over 5 seeds (the
selection metric); \textbf{ens.\ AUC} is the 5-seed probability ensemble. Adopted
Stage-2 finalists in \textbf{bold}.}
\label{tab:ft-grid}
\centering
\small
\setlength{\tabcolsep}{4pt}
\begin{tabular}{@{}lcc@{}}
\toprule
\textbf{Configuration} & \textbf{AUC (per-seed)} & \textbf{ens.\ AUC} \\
\midrule
\textbf{mxbai, LoRA $r$16, lr1e-4}  & \textbf{0.916}{\footnotesize$\pm$.005} & 0.917 \\
\textbf{LoRA $r$16+FFN, lr1e-4}     & \textbf{0.915}{\footnotesize$\pm$.008} & 0.920 \\
\textbf{LoRA $r$8, lr2e-4}          & \textbf{0.913}{\footnotesize$\pm$.010} & 0.924 \\
\textbf{last-4, lr1e-5}             & \textbf{0.908}{\footnotesize$\pm$.008} & 0.899 \\
LoRA $r$16, lr2e-4                  & 0.907{\footnotesize$\pm$.022} & 0.920 \\
full, lr1e-5                        & 0.906{\footnotesize$\pm$.006} & 0.899 \\
last-2, lr2e-5                      & 0.899{\footnotesize$\pm$.005} & 0.899 \\
BitFit, lr1e-4                      & 0.899{\footnotesize$\pm$.010} & 0.902 \\
LoRA $r$32, lr1e-4                  & 0.893{\footnotesize$\pm$.053} & 0.924 \\
frozen (control)                   & 0.892{\footnotesize$\pm$.003} & 0.884 \\
mxbai, frozen                      & 0.886{\footnotesize$\pm$.003} & 0.884 \\
LoRA $r$16, lr1e-4                  & 0.888{\footnotesize$\pm$.052} & 0.920 \\
\midrule
DAPT, frozen                       & 0.709{\footnotesize$\pm$.044} & 0.707 \\
DAPT, LoRA $r$16                    & 0.708{\footnotesize$\pm$.078} & 0.649 \\
\bottomrule
\end{tabular}
\end{table}

Three findings emerge. \textbf{(1) Encoder adaptation gives a modest, real
lift.} The best single-model configurations reach a per-seed mean AUC of
$\approx\!0.915$ against the frozen control's $0.892$ --- a $\approx\!+0.02$ AUC
gain that clears the pre-registered adoption band. The lift is smaller and the
field tighter than the ensemble metric suggests: the top LoRA, \texttt{last-}4,
and full configurations ($0.906$--$0.916$) are all within seed noise of one
another, so it does \emph{not} establish a single method as the winner
--- only that lightweight adaptation helps and that the choice among
parameter-efficient methods must be settled on the test split. We nonetheless
note that LoRA supplies the most \emph{stable} top configurations (low per-seed
variance at high mean), whereas the configurations the ensemble metric would
have promoted --- LoRA $r$32/$r$16 at lr1e-4 --- carry $\pm.05$ seed variance and
sit at single-model rank $\approx\!11$, a cautionary example of the
ensemble--single-model mismatch. \textbf{(2) The encoder swap is neutral.}
\texttt{mxbai} matches \texttt{bge} under both frozen ($0.886$ vs.\ $0.892$) and
LoRA ($0.916$ vs.\ $0.915$); it is the nominal rank-1 single model but is within
noise of the \texttt{bge} leaders. \textbf{(3) DAPT collapses.} Both DAPT arms
fall to $0.65$--$0.71$ AUC --- far \emph{below} the no-adaptation baseline ---
and predict the majority class (per-seed F1 $=0$ on several seeds). The
masked-LM objective, applied to a 107-interview corpus with no
embedding-preserving term, degrades the very sentence-embedding geometry the head
pools over: a catastrophic-forgetting failure, not an under-training one.
Lengthening the masked-LM schedule would be expected to worsen it; reviving DAPT
would need an embedding-preserving objective (e.g.\ TSDAE or contrastive
continued pre-training), which we do not pursue.

The pre-registered rule promotes the four configurations within the seed-noise
band of the best (capped at four) --- \texttt{mxbai} LoRA $r$16,
\texttt{bge} LoRA $r$16+FFN, \texttt{bge} LoRA $r$8/lr2e-4, and \texttt{last-}4 ---
plus the frozen control to Stage 2, and carries the rank-1 \texttt{mxbai} LoRA
$r$16 forward as the single winner for the Stage-3 cross-validation.

\subsection{The leakage-free OOF probe}
\label{sec:finetune-iter1-oof}

Before the once-only test evaluation, each finalist passes through the Stage-2
out-of-fold (OOF) probe: a nested $K$-fold over the train+dev pool that gives
every one of the $N\!=\!142$ subjects a leakage-free probability from a model
that never trained on it. The probe exists to fix the decision threshold, but it
is \emph{also} a leakage-free, fourfold-larger-$N$ read on the finalist ranking
than the $n\!=\!35$ development split, and it is informative on its own.
\Cref{tab:ft-oof} reports it.

\begin{table}[t]
\caption{Stage-2 OOF probe (leakage-free, train+dev pool $N\!=\!142$, 42
positive). \textbf{OOF AUC} with a subject-resampling bootstrap 95\% CI; the
development and OOF ranks are repeated for contrast (development per-seed AUC in
\Cref{tab:ft-grid}).}
\label{tab:ft-oof}
\centering\small
\setlength{\tabcolsep}{4pt}
\begin{tabular}{@{}lcccc@{}}
\toprule
\textbf{Model} & \textbf{OOF AUC} & \textbf{95\% CI} & \textbf{dev rk} & \textbf{OOF rk} \\
\midrule
\textbf{frozen (control)}      & \textbf{0.790} & [.709,.866] & 5 & \textbf{1} \\
last-4                         & 0.789 & [.706,.862] & 4 & 2 \\
LoRA $r$16+FFN                 & 0.775 & [.687,.851] & 2 & 3 \\
mxbai LoRA $r$16               & 0.768 & [.685,.843] & \textbf{1} & 4 \\
LoRA $r$8                      & 0.764 & [.676,.846] & 3 & 5 \\
\bottomrule
\end{tabular}
\end{table}

Three points follow, all consistent with the development-grid reading.
\textbf{(1) The grid ranking does not survive the leakage-free probe.} The
development rank-1 configuration (\texttt{mxbai} LoRA $r$16) falls to OOF rank~4,
while \texttt{last-}4 --- development rank~4 --- rises to OOF rank~2; yet every
pairwise OOF AUC gap lies within bootstrap noise, so the finalists are
\emph{statistically indistinguishable} at $N\!=\!142$. This is the same
conclusion the per-seed grid reached --- no single adaptation method is crowned
--- now corroborated on a leakage-free, larger sample, and it is
exactly why selection was pre-registered and the test split is spent once.
\textbf{(2) The development AUC is optimistic.} The leakage-free OOF AUC
($\approx\!0.76$--$0.79$) sits about $0.13$ below the development figure
($\approx\!0.91$); test expectations should be calibrated to the OOF band.
\textbf{(3) Fine-tuning yields no leakage-free OOF lift over the frozen
control.} Run through the identical OOF probe, the frozen encoder scores the
\emph{highest} pooled OOF AUC of the set ($0.790$, $95\%$ CI $[.709,.866]$),
edging the best adaptation arm (\texttt{last-}4, $0.789$). Treating the 15 paired
per-run AUCs ($5$ fold $\times\,3$ seed, aligned by fold and seed) with a
Wilcoxon signed-rank test, no fine-tuning finalist differs from frozen: the
mean per-run gaps are within $\pm0.01$ AUC and the raw $p$-values are
$0.98$/$0.39$/$0.60$/$0.69$ (\texttt{last-}4/LoRA $r$16+FFN/\texttt{mxbai}/LoRA
$r$8), all $\approx\!1$ after Holm correction. Every model separates the pooled
pool far above chance (Mann-Whitney $p<10^{-6}$), but adaptation moves the
\emph{ranking} no further than frozen does. The grid's apparent $\approx\!+0.02$
development-AUC lift over the frozen control is therefore a small-$n$ artefact of
the $n\!=\!35$ dev split; it does not reproduce on the four-fold-larger,
leakage-free OOF read. This is exactly the kind of optimism the staged protocol
was built to catch before spending the test split.

\subsection{Official-test and cross-validation results}
\label{sec:finetune-iter1-test}

The leakage-free OOF probe predicted that the development lift would not survive.
The two once-only readings --- the official AVEC2017 test split and the
cross-validation pool --- confirm it.

\paragraph{Official test split.} Each finalist is evaluated once on the official
test split ($n\!=\!47$, 14 positive) as a 30-seed probability ensemble, with the
decision threshold fixed by the leakage-free OOF prevalence rule of
\Cref{sec:finetune-iter1-oof}; \Cref{tab:ft-test} reports the result against the
frozen single-model bar (AUC $0.864$, macro-F1 $0.774$; lever~3 of
\Cref{tab:levers}). \emph{No fine-tuning finalist beats the frozen bar.} The best
ROC-AUC --- \texttt{last-}4 at $0.857$ --- still sits below frozen, and the
remaining three span $0.816$--$0.846$; every macro-F1 ($0.743$--$0.763$) is below
the frozen $0.774$. The comparison is in fact conservative against frozen: these
are 30-seed \emph{ensemble} predictions, yet they fail to clear a \emph{single}
frozen model. Each AUC bootstrap CI is wide ($\approx\!0.24$ span at $n\!=\!47$)
and contains the frozen bar, so no finalist is \emph{distinguishable} from frozen
either way --- the honest reading is a null, not a regression. The
development-grid ordering again fails to transfer: the dev rank-1 model
(\texttt{mxbai} LoRA $r$16) posts the \emph{lowest} test AUC ($0.816$), while the
dev rank-4 \texttt{last-}4 posts the highest --- the same rank inversion the OOF
probe showed, now reproduced on the test split.

\begin{table}[t]
\caption{Official AVEC2017 test split ($n\!=\!47$, 14 positive), 30-seed
probability ensemble at the leakage-free OOF prevalence-matched threshold.
\textbf{AUC} carries a subject-resampling bootstrap 95\% CI. Top row is the frozen
single-model bar (AUC $0.864$, macro-F1 $0.774$). Rows ordered by test AUC; the
development rank is repeated for contrast.}
\label{tab:ft-test}
\centering\small
\setlength{\tabcolsep}{4pt}
\begin{tabular}{@{}lcccc@{}}
\toprule
\textbf{Model} & \textbf{AUC} & \textbf{95\% CI} & \textbf{mac-F1} & \textbf{dev rk} \\
\midrule
frozen (bar)                   & \textbf{0.864} & --- & \textbf{0.774} & --- \\
last-4, lr1e-5                 & 0.857 & [.722,.964] & 0.743 & 4 \\
LoRA $r$16+FFN, lr1e-4         & 0.846 & [.706,.959] & 0.743 & 2 \\
LoRA $r$8, lr2e-4              & 0.844 & [.708,.953] & 0.743 & 3 \\
mxbai LoRA $r$16, lr1e-4       & 0.816 & [.676,.933] & 0.763 & 1 \\
\bottomrule
\end{tabular}
\end{table}

\paragraph{Cross-validation.} Stage~3 re-evaluates the single winner
(\texttt{mxbai} LoRA $r$16, carried forward as the nominal dev rank-1) against the
frozen control under repeated $K$-fold cross-validation on the train+dev pool, all
adaptation performed inside each fold (\Cref{tab:ft-cv}). Over the 50 paired
per-run AUCs (2 repeats $\times\,5$ fold $\times\,5$ seed, aligned by repeat, fold,
and seed), the winner and frozen are indistinguishable on AUC ($0.777$ vs.\ $0.778$;
Wilcoxon $p=0.79$) and frozen is nominally ahead on macro-F1 ($0.679$ vs.\ $0.657$;
$p=0.12$, n.s.). Monte-Carlo cross-validation, run for both arms,
\emph{flips the sign} of the nominal gap --- the winner
edges frozen on MC (per-run AUC $0.824$ vs.\ $0.803$, $\Delta\!=\!+0.021$,
$p=0.35$; macro-F1 $0.694$ vs.\ $0.667$, $\Delta\!=\!+0.026$, $p=0.47$) having
trailed it on $K$-fold --- but neither gap is significant, and a lift that
reverses direction between two cross-validation schemes is noise, not signal.

\begin{table}[t]
\caption{Stage-3 cross-validation (train+dev pool), winner (\texttt{mxbai} LoRA
$r$16) vs.\ the in-harness frozen control, adaptation inside each fold. Repeated
$K$-fold is $2$ repeat $\times\,5$ fold $\times\,5$ seed $=50$ per-run
evaluations; Monte-Carlo is $25$ per-run. Per-run mean$\pm$std and the
fold-ensemble.}
\label{tab:ft-cv}
\centering\small
\setlength{\tabcolsep}{4pt}
\begin{tabular}{@{}llcccc@{}}
\toprule
 & & \multicolumn{2}{c}{\textbf{per-run AUC}} & \multicolumn{2}{c}{\textbf{per-run mac-F1}} \\
\cmidrule(lr){3-4}\cmidrule(lr){5-6}
\textbf{Scheme} & \textbf{Model} & \textbf{run} & \textbf{ens.} & \textbf{run} & \textbf{ens.} \\
\midrule
$K$-fold & frozen        & \textbf{0.778}{\footnotesize$\pm$.11} & 0.80 & \textbf{0.679}{\footnotesize$\pm$.14} & 0.72 \\
$K$-fold & winner        & 0.777{\footnotesize$\pm$.11} & 0.80 & 0.657{\footnotesize$\pm$.11} & 0.67 \\
\midrule
MC       & frozen        & 0.803{\footnotesize$\pm$.09} & 0.843 & 0.667{\footnotesize$\pm$.12} & 0.732 \\
MC       & winner        & \textbf{0.824}{\footnotesize$\pm$.09} & \textbf{0.847} & \textbf{0.694}{\footnotesize$\pm$.13} & \textbf{0.749} \\
\bottomrule
\end{tabular}
\end{table}

\paragraph{Verdict.} Three independent leakage-free readings --- the
$N\!=\!142$ OOF probe, the once-only $n\!=\!47$ test split, and the repeated
cross-validation pool --- agree: encoder fine-tuning yields \emph{no} lift over the
frozen baseline for this task at this scale. The $\approx\!+0.02$ AUC advantage on
the $n\!=\!35$ development grid was a small-sample artefact; it vanishes the moment
it is measured on a larger, leakage-free sample, and the dev ranking among
adaptation methods never transfers (the dev rank-1 \texttt{mxbai} LoRA $r$16 is the
worst finalist at test). The pre-registered, staged protocol did exactly what it
was designed to do: it exposed the development optimism \emph{before} the test
split was spent, and prevents us from over-claiming a parameter-efficient
fine-tuning win that the data does not support. We therefore retain the frozen
encoder as the deployed system; the remaining headroom (\Cref{tab:levers}) is not
recoverable by fine-tuning the sentence encoder on a 107-interview corpus.

% placed AFTER \section{Internal Comparison: Cross-Validation}
% (\label{sec:internal}) and BEFORE \section{Discussion and Limitations}.
% Iteration 1 is complete on the development grid; the official-test and
% cross-validation blocks marked %TODO are filled once cluster jobs 3/4 land.
% =====================================================================
\section{Does Fine-Tuning the Encoder Help?}
\label{sec:finetune}

Every result so far keeps the sentence encoder \emph{frozen}: the
single-lever study (\Cref{tab:levers}) showed the remaining headroom is in the
\emph{ranking}, not the decision rule, therefore to lift the AUC ceiling a possibility is
to fine-tune the encoder.
Fine-tuning an $n\!=\!107$ corpus is the regime where over-fitting is most
likely, so the study is run as a pre-registered, staged search that selects on
development data and touches the official test split only once, for a small set
of finalists.

\subsection{Fine-tuning menu and staged protocol}
\label{sec:finetune-protocol}
We compare five adaptation regimes against the frozen baseline, ordered by
trainable-parameter budget: (i) \textbf{frozen} (control, embeddings only);
(ii) \textbf{BitFit} (bias terms only); (iii) \textbf{LoRA} low-rank adapters on
the attention (optionally FFN) projections, rank $r\!\in\!\{8,16,32\}$;
(iv) \textbf{last-$k$} unfreezing of the top $k\!\in\!\{2,4\}$ transformer
blocks with layer-wise LR decay; and (v) \textbf{full} fine-tuning. We also test
\textbf{domain-adaptive pre-training} (DAPT): a masked-LM pass over the training
interviews before the supervised stage. The encoder is \texttt{bge-large-en-v1.5}
throughout; one arm swaps in \texttt{mxbai-embed-large-v1} to check
encoder sensitivity. The temporal GRU pooling head
(\Cref{sec:method}) is unchanged.

Selection is staged to ring-fence the test split:
\begin{itemize}
  \item \textbf{Stage 1 --- development grid.} 20 configurations, 5 seeds each,
        \emph{official protocol but no test evaluation}. Configurations are ranked on the 5-seed
        \emph{per-seed mean} development ROC-AUC --- the same quantity reported
        at test time.
  \item \textbf{Stage 2 --- finalists.} Every configuration whose per-seed mean
        development AUC is within one seed-noise band ($\pm0.015$) of the best,
        capped at four, plus the frozen control, advances. Each gets a
        leakage-free decision threshold from a nested out-of-fold (OOF) probe and
        is then evaluated \emph{once} on the official test split over 30 seeds.
  \item \textbf{Stage 3 --- cross-validation.} The single winner and the frozen
        control are re-evaluated under $K$-fold (with a repeat) and
        Monte-Carlo cross-validation on the training pool, mirroring
        \Cref{sec:internal}.
\end{itemize}
All encoder adaptation happens \emph{inside} each fold; no test or
held-out interview is ever seen during fine-tuning or thresholding.

\subsection{Development-grid selection}
\label{sec:finetune-iter1}

A practical caveat governs how the grid is read. The development split is
$n\!=\!35$ interviews (12 positive); at this size the \emph{thresholded} metrics
are quantised --- many unrelated configurations land on the identical confusion
matrix --- so they carry little ranking signal. We therefore rank on the
per-seed mean ROC-AUC, the only continuous discriminator, and treat AUC gaps
below $\approx\!0.01$ as noise. \Cref{tab:ft-grid} reports the representative
configurations.

These dev figures are produced by the fine-tuning harness --- which trains the
MIL head with a shorter, more strongly regularised schedule than the
design-ablation pipeline of \Cref{tab:abldesign-rep} (and embeds under mixed
precision) --- so they are internally consistent but \emph{not} numerically
comparable to that table's frozen dev AUC. All fine-tuning effects are therefore
measured against the in-harness frozen control, not against the design-ablation
numbers; the headline frozen reference remains the once-only official-test bar.

\begin{table}[t]
\caption{Fine-tuning grid, development split, ranked by the
selection metric. \textbf{AUC} is the per-seed mean$\pm$std over 5 seeds (the
selection metric); \textbf{ens.\ AUC} is the 5-seed probability ensemble. Adopted
Stage-2 finalists in \textbf{bold}.}
\label{tab:ft-grid}
\centering
\small
\setlength{\tabcolsep}{4pt}
\begin{tabular}{@{}lcc@{}}
\toprule
\textbf{Configuration} & \textbf{AUC (per-seed)} & \textbf{ens.\ AUC} \\
\midrule
\textbf{mxbai, LoRA $r$16, lr1e-4}  & \textbf{0.916}{\footnotesize$\pm$.005} & 0.917 \\
\textbf{LoRA $r$16+FFN, lr1e-4}     & \textbf{0.915}{\footnotesize$\pm$.008} & 0.920 \\
\textbf{LoRA $r$8, lr2e-4}          & \textbf{0.913}{\footnotesize$\pm$.010} & 0.924 \\
\textbf{last-4, lr1e-5}             & \textbf{0.908}{\footnotesize$\pm$.008} & 0.899 \\
LoRA $r$16, lr2e-4                  & 0.907{\footnotesize$\pm$.022} & 0.920 \\
full, lr1e-5                        & 0.906{\footnotesize$\pm$.006} & 0.899 \\
last-2, lr2e-5                      & 0.899{\footnotesize$\pm$.005} & 0.899 \\
BitFit, lr1e-4                      & 0.899{\footnotesize$\pm$.010} & 0.902 \\
LoRA $r$32, lr1e-4                  & 0.893{\footnotesize$\pm$.053} & 0.924 \\
frozen (control)                   & 0.892{\footnotesize$\pm$.003} & 0.884 \\
mxbai, frozen                      & 0.886{\footnotesize$\pm$.003} & 0.884 \\
LoRA $r$16, lr1e-4                  & 0.888{\footnotesize$\pm$.052} & 0.920 \\
\midrule
DAPT, frozen                       & 0.709{\footnotesize$\pm$.044} & 0.707 \\
DAPT, LoRA $r$16                    & 0.708{\footnotesize$\pm$.078} & 0.649 \\
\bottomrule
\end{tabular}
\end{table}

Three findings emerge. \textbf{(1) Encoder adaptation gives a modest, real
lift.} The best single-model configurations reach a per-seed mean AUC of
$\approx\!0.915$ against the frozen control's $0.892$ --- a $\approx\!+0.02$ AUC
gain that clears the pre-registered adoption band. The lift is smaller and the
field tighter than the ensemble metric suggests: the top LoRA, \texttt{last-}4,
and full configurations ($0.906$--$0.916$) are all within seed noise of one
another, so it does \emph{not} establish a single method as the winner
--- only that lightweight adaptation helps and that the choice among
parameter-efficient methods must be settled on the test split. We nonetheless
note that LoRA supplies the most \emph{stable} top configurations (low per-seed
variance at high mean), whereas the configurations the ensemble metric would
have promoted --- LoRA $r$32/$r$16 at lr1e-4 --- carry $\pm.05$ seed variance and
sit at single-model rank $\approx\!11$, a cautionary example of the
ensemble--single-model mismatch. \textbf{(2) The encoder swap is neutral.}
\texttt{mxbai} matches \texttt{bge} under both frozen ($0.886$ vs.\ $0.892$) and
LoRA ($0.916$ vs.\ $0.915$); it is the nominal rank-1 single model but is within
noise of the \texttt{bge} leaders. \textbf{(3) DAPT collapses.} Both DAPT arms
fall to $0.65$--$0.71$ AUC --- far \emph{below} the no-adaptation baseline ---
and predict the majority class (per-seed F1 $=0$ on several seeds). The
masked-LM objective, applied to a 107-interview corpus with no
embedding-preserving term, degrades the very sentence-embedding geometry the head
pools over: a catastrophic-forgetting failure, not an under-training one.
Lengthening the masked-LM schedule would be expected to worsen it; reviving DAPT
would need an embedding-preserving objective (e.g.\ TSDAE or contrastive
continued pre-training), which we do not pursue.

The pre-registered rule promotes the four configurations within the seed-noise
band of the best (capped at four) --- \texttt{mxbai} LoRA $r$16,
\texttt{bge} LoRA $r$16+FFN, \texttt{bge} LoRA $r$8/lr2e-4, and \texttt{last-}4 ---
plus the frozen control to Stage 2, and carries the rank-1 \texttt{mxbai} LoRA
$r$16 forward as the single winner for the Stage-3 cross-validation.

\subsection{The leakage-free OOF probe}
\label{sec:finetune-iter1-oof}

Before the once-only test evaluation, each finalist passes through the Stage-2
out-of-fold (OOF) probe: a nested $K$-fold over the train+dev pool that gives
every one of the $N\!=\!142$ subjects a leakage-free probability from a model
that never trained on it. The probe exists to fix the decision threshold, but it
is \emph{also} a leakage-free, fourfold-larger-$N$ read on the finalist ranking
than the $n\!=\!35$ development split, and it is informative on its own.
\Cref{tab:ft-oof} reports it.

\begin{table}[t]
\caption{Stage-2 OOF probe (leakage-free, train+dev pool $N\!=\!142$, 42
positive). \textbf{OOF AUC} with a subject-resampling bootstrap 95\% CI; the
development and OOF ranks are repeated for contrast (development per-seed AUC in
\Cref{tab:ft-grid}).}
\label{tab:ft-oof}
\centering\small
\setlength{\tabcolsep}{4pt}
\begin{tabular}{@{}lcccc@{}}
\toprule
\textbf{Model} & \textbf{OOF AUC} & \textbf{95\% CI} & \textbf{dev rk} & \textbf{OOF rk} \\
\midrule
\textbf{frozen (control)}      & \textbf{0.790} & [.709,.866] & 5 & \textbf{1} \\
last-4                         & 0.789 & [.706,.862] & 4 & 2 \\
LoRA $r$16+FFN                 & 0.775 & [.687,.851] & 2 & 3 \\
mxbai LoRA $r$16               & 0.768 & [.685,.843] & \textbf{1} & 4 \\
LoRA $r$8                      & 0.764 & [.676,.846] & 3 & 5 \\
\bottomrule
\end{tabular}
\end{table}

Three points follow, all consistent with the development-grid reading.
\textbf{(1) The grid ranking does not survive the leakage-free probe.} The
development rank-1 configuration (\texttt{mxbai} LoRA $r$16) falls to OOF rank~4,
while \texttt{last-}4 --- development rank~4 --- rises to OOF rank~2; yet every
pairwise OOF AUC gap lies within bootstrap noise, so the finalists are
\emph{statistically indistinguishable} at $N\!=\!142$. This is the same
conclusion the per-seed grid reached --- no single adaptation method is crowned
--- now corroborated on a leakage-free, larger sample, and it is
exactly why selection was pre-registered and the test split is spent once.
\textbf{(2) The development AUC is optimistic.} The leakage-free OOF AUC
($\approx\!0.76$--$0.79$) sits about $0.13$ below the development figure
($\approx\!0.91$); test expectations should be calibrated to the OOF band.
\textbf{(3) Fine-tuning yields no leakage-free OOF lift over the frozen
control.} Run through the identical OOF probe, the frozen encoder scores the
\emph{highest} pooled OOF AUC of the set ($0.790$, $95\%$ CI $[.709,.866]$),
edging the best adaptation arm (\texttt{last-}4, $0.789$). Treating the 15 paired
per-run AUCs ($5$ fold $\times\,3$ seed, aligned by fold and seed) with a
Wilcoxon signed-rank test, no fine-tuning finalist differs from frozen: the
mean per-run gaps are within $\pm0.01$ AUC and the raw $p$-values are
$0.98$/$0.39$/$0.60$/$0.69$ (\texttt{last-}4/LoRA $r$16+FFN/\texttt{mxbai}/LoRA
$r$8), all $\approx\!1$ after Holm correction. Every model separates the pooled
pool far above chance (Mann-Whitney $p<10^{-6}$), but adaptation moves the
\emph{ranking} no further than frozen does. The grid's apparent $\approx\!+0.02$
development-AUC lift over the frozen control is therefore a small-$n$ artefact of
the $n\!=\!35$ dev split; it does not reproduce on the four-fold-larger,
leakage-free OOF read. This is exactly the kind of optimism the staged protocol
was built to catch before spending the test split.

\subsection{Official-test and cross-validation results}
\label{sec:finetune-iter1-test}

The leakage-free OOF probe predicted that the development lift would not survive.
The two once-only readings --- the official AVEC2017 test split and the
cross-validation pool --- confirm it.

\paragraph{Official test split.} Each finalist is evaluated once on the official
test split ($n\!=\!47$, 14 positive) as a 30-seed probability ensemble, with the
decision threshold fixed by the leakage-free OOF prevalence rule of
\Cref{sec:finetune-iter1-oof}; \Cref{tab:ft-test} reports the result against the
frozen single-model bar (AUC $0.864$, macro-F1 $0.774$; lever~3 of
\Cref{tab:levers}). \emph{No fine-tuning finalist beats the frozen bar.} The best
ROC-AUC --- \texttt{last-}4 at $0.857$ --- still sits below frozen, and the
remaining three span $0.816$--$0.846$; every macro-F1 ($0.743$--$0.763$) is below
the frozen $0.774$. The comparison is in fact conservative against frozen: these
are 30-seed \emph{ensemble} predictions, yet they fail to clear a \emph{single}
frozen model. Each AUC bootstrap CI is wide ($\approx\!0.24$ span at $n\!=\!47$)
and contains the frozen bar, so no finalist is \emph{distinguishable} from frozen
either way --- the honest reading is a null, not a regression. The
development-grid ordering again fails to transfer: the dev rank-1 model
(\texttt{mxbai} LoRA $r$16) posts the \emph{lowest} test AUC ($0.816$), while the
dev rank-4 \texttt{last-}4 posts the highest --- the same rank inversion the OOF
probe showed, now reproduced on the test split.

\begin{table}[t]
\caption{Official AVEC2017 test split ($n\!=\!47$, 14 positive), 30-seed
probability ensemble at the leakage-free OOF prevalence-matched threshold.
\textbf{AUC} carries a subject-resampling bootstrap 95\% CI. Top row is the frozen
single-model bar (AUC $0.864$, macro-F1 $0.774$). Rows ordered by test AUC; the
development rank is repeated for contrast.}
\label{tab:ft-test}
\centering\small
\setlength{\tabcolsep}{4pt}
\begin{tabular}{@{}lcccc@{}}
\toprule
\textbf{Model} & \textbf{AUC} & \textbf{95\% CI} & \textbf{mac-F1} & \textbf{dev rk} \\
\midrule
frozen (bar)                   & \textbf{0.864} & --- & \textbf{0.774} & --- \\
last-4, lr1e-5                 & 0.857 & [.722,.964] & 0.743 & 4 \\
LoRA $r$16+FFN, lr1e-4         & 0.846 & [.706,.959] & 0.743 & 2 \\
LoRA $r$8, lr2e-4              & 0.844 & [.708,.953] & 0.743 & 3 \\
mxbai LoRA $r$16, lr1e-4       & 0.816 & [.676,.933] & 0.763 & 1 \\
\bottomrule
\end{tabular}
\end{table}

\paragraph{Cross-validation.} Stage~3 re-evaluates the single winner
(\texttt{mxbai} LoRA $r$16, carried forward as the nominal dev rank-1) against the
frozen control under repeated $K$-fold cross-validation on the train+dev pool, all
adaptation performed inside each fold (\Cref{tab:ft-cv}). Over the 50 paired
per-run AUCs (2 repeats $\times\,5$ fold $\times\,5$ seed, aligned by repeat, fold,
and seed), the winner and frozen are indistinguishable on AUC ($0.777$ vs.\ $0.778$;
Wilcoxon $p=0.79$) and frozen is nominally ahead on macro-F1 ($0.679$ vs.\ $0.657$;
$p=0.12$, n.s.). Monte-Carlo cross-validation, run for both arms,
\emph{flips the sign} of the nominal gap --- the winner
edges frozen on MC (per-run AUC $0.824$ vs.\ $0.803$, $\Delta\!=\!+0.021$,
$p=0.35$; macro-F1 $0.694$ vs.\ $0.667$, $\Delta\!=\!+0.026$, $p=0.47$) having
trailed it on $K$-fold --- but neither gap is significant, and a lift that
reverses direction between two cross-validation schemes is noise, not signal.

\begin{table}[t]
\caption{Stage-3 cross-validation (train+dev pool), winner (\texttt{mxbai} LoRA
$r$16) vs.\ the in-harness frozen control, adaptation inside each fold. Repeated
$K$-fold is $2$ repeat $\times\,5$ fold $\times\,5$ seed $=50$ per-run
evaluations; Monte-Carlo is $25$ per-run. Per-run mean$\pm$std and the
fold-ensemble.}
\label{tab:ft-cv}
\centering\small
\setlength{\tabcolsep}{4pt}
\begin{tabular}{@{}llcccc@{}}
\toprule
 & & \multicolumn{2}{c}{\textbf{per-run AUC}} & \multicolumn{2}{c}{\textbf{per-run mac-F1}} \\
\cmidrule(lr){3-4}\cmidrule(lr){5-6}
\textbf{Scheme} & \textbf{Model} & \textbf{run} & \textbf{ens.} & \textbf{run} & \textbf{ens.} \\
\midrule
$K$-fold & frozen        & \textbf{0.778}{\footnotesize$\pm$.11} & 0.80 & \textbf{0.679}{\footnotesize$\pm$.14} & 0.72 \\
$K$-fold & winner        & 0.777{\footnotesize$\pm$.11} & 0.80 & 0.657{\footnotesize$\pm$.11} & 0.67 \\
\midrule
MC       & frozen        & 0.803{\footnotesize$\pm$.09} & 0.843 & 0.667{\footnotesize$\pm$.12} & 0.732 \\
MC       & winner        & \textbf{0.824}{\footnotesize$\pm$.09} & \textbf{0.847} & \textbf{0.694}{\footnotesize$\pm$.13} & \textbf{0.749} \\
\bottomrule
\end{tabular}
\end{table}

\paragraph{Verdict.} Three independent leakage-free readings --- the
$N\!=\!142$ OOF probe, the once-only $n\!=\!47$ test split, and the repeated
cross-validation pool --- agree: encoder fine-tuning yields \emph{no} lift over the
frozen baseline for this task at this scale. The $\approx\!+0.02$ AUC advantage on
the $n\!=\!35$ development grid was a small-sample artefact; it vanishes the moment
it is measured on a larger, leakage-free sample, and the dev ranking among
adaptation methods never transfers (the dev rank-1 \texttt{mxbai} LoRA $r$16 is the
worst finalist at test). The pre-registered, staged protocol did exactly what it
was designed to do: it exposed the development optimism \emph{before} the test
split was spent, and prevents us from over-claiming a parameter-efficient
fine-tuning win that the data does not support. We therefore retain the frozen
encoder as the deployed system; the remaining headroom (\Cref{tab:levers}) is not
recoverable by fine-tuning the sentence encoder on a 107-interview corpus.

% placed AFTER \section{Internal Comparison: Cross-Validation}
% (\label{sec:internal}) and BEFORE \section{Discussion and Limitations}.
% Iteration 1 is complete on the development grid; the official-test and
% cross-validation blocks marked %TODO are filled once cluster jobs 3/4 land.
% =====================================================================
\section{Does Fine-Tuning the Encoder Help?}
\label{sec:finetune}

Every result so far keeps the sentence encoder \emph{frozen}: the
single-lever study (\Cref{tab:levers}) showed the remaining headroom is in the
\emph{ranking}, not the decision rule, therefore to lift the AUC ceiling a possibility is
to fine-tune the encoder.
Fine-tuning an $n\!=\!107$ corpus is the regime where over-fitting is most
likely, so the study is run as a pre-registered, staged search that selects on
development data and touches the official test split only once, for a small set
of finalists.

\subsection{Fine-tuning menu and staged protocol}
\label{sec:finetune-protocol}
We compare five adaptation regimes against the frozen baseline, ordered by
trainable-parameter budget: (i) \textbf{frozen} (control, embeddings only);
(ii) \textbf{BitFit} (bias terms only); (iii) \textbf{LoRA} low-rank adapters on
the attention (optionally FFN) projections, rank $r\!\in\!\{8,16,32\}$;
(iv) \textbf{last-$k$} unfreezing of the top $k\!\in\!\{2,4\}$ transformer
blocks with layer-wise LR decay; and (v) \textbf{full} fine-tuning. We also test
\textbf{domain-adaptive pre-training} (DAPT): a masked-LM pass over the training
interviews before the supervised stage. The encoder is \texttt{bge-large-en-v1.5}
throughout; one arm swaps in \texttt{mxbai-embed-large-v1} to check
encoder sensitivity. The temporal GRU pooling head
(\Cref{sec:method}) is unchanged.

Selection is staged to ring-fence the test split:
\begin{itemize}
  \item \textbf{Stage 1 --- development grid.} 20 configurations, 5 seeds each,
        \emph{official protocol but no test evaluation}. Configurations are ranked on the 5-seed
        \emph{per-seed mean} development ROC-AUC --- the same quantity reported
        at test time.
  \item \textbf{Stage 2 --- finalists.} Every configuration whose per-seed mean
        development AUC is within one seed-noise band ($\pm0.015$) of the best,
        capped at four, plus the frozen control, advances. Each gets a
        leakage-free decision threshold from a nested out-of-fold (OOF) probe and
        is then evaluated \emph{once} on the official test split over 30 seeds.
  \item \textbf{Stage 3 --- cross-validation.} The single winner and the frozen
        control are re-evaluated under $K$-fold (with a repeat) and
        Monte-Carlo cross-validation on the training pool, mirroring
        \Cref{sec:internal}.
\end{itemize}
All encoder adaptation happens \emph{inside} each fold; no test or
held-out interview is ever seen during fine-tuning or thresholding.

\subsection{Development-grid selection}
\label{sec:finetune-iter1}

A practical caveat governs how the grid is read. The development split is
$n\!=\!35$ interviews (12 positive); at this size the \emph{thresholded} metrics
are quantised --- many unrelated configurations land on the identical confusion
matrix --- so they carry little ranking signal. We therefore rank on the
per-seed mean ROC-AUC, the only continuous discriminator, and treat AUC gaps
below $\approx\!0.01$ as noise. \Cref{tab:ft-grid} reports the representative
configurations.

These dev figures are produced by the fine-tuning harness --- which trains the
MIL head with a shorter, more strongly regularised schedule than the
design-ablation pipeline of \Cref{tab:abldesign-rep} (and embeds under mixed
precision) --- so they are internally consistent but \emph{not} numerically
comparable to that table's frozen dev AUC. All fine-tuning effects are therefore
measured against the in-harness frozen control, not against the design-ablation
numbers; the headline frozen reference remains the once-only official-test bar.

\begin{table}[t]
\caption{Fine-tuning grid, development split, ranked by the
selection metric. \textbf{AUC} is the per-seed mean$\pm$std over 5 seeds (the
selection metric); \textbf{ens.\ AUC} is the 5-seed probability ensemble. Adopted
Stage-2 finalists in \textbf{bold}.}
\label{tab:ft-grid}
\centering
\small
\setlength{\tabcolsep}{4pt}
\begin{tabular}{@{}lcc@{}}
\toprule
\textbf{Configuration} & \textbf{AUC (per-seed)} & \textbf{ens.\ AUC} \\
\midrule
\textbf{mxbai, LoRA $r$16, lr1e-4}  & \textbf{0.916}{\footnotesize$\pm$.005} & 0.917 \\
\textbf{LoRA $r$16+FFN, lr1e-4}     & \textbf{0.915}{\footnotesize$\pm$.008} & 0.920 \\
\textbf{LoRA $r$8, lr2e-4}          & \textbf{0.913}{\footnotesize$\pm$.010} & 0.924 \\
\textbf{last-4, lr1e-5}             & \textbf{0.908}{\footnotesize$\pm$.008} & 0.899 \\
LoRA $r$16, lr2e-4                  & 0.907{\footnotesize$\pm$.022} & 0.920 \\
full, lr1e-5                        & 0.906{\footnotesize$\pm$.006} & 0.899 \\
last-2, lr2e-5                      & 0.899{\footnotesize$\pm$.005} & 0.899 \\
BitFit, lr1e-4                      & 0.899{\footnotesize$\pm$.010} & 0.902 \\
LoRA $r$32, lr1e-4                  & 0.893{\footnotesize$\pm$.053} & 0.924 \\
frozen (control)                   & 0.892{\footnotesize$\pm$.003} & 0.884 \\
mxbai, frozen                      & 0.886{\footnotesize$\pm$.003} & 0.884 \\
LoRA $r$16, lr1e-4                  & 0.888{\footnotesize$\pm$.052} & 0.920 \\
\midrule
DAPT, frozen                       & 0.709{\footnotesize$\pm$.044} & 0.707 \\
DAPT, LoRA $r$16                    & 0.708{\footnotesize$\pm$.078} & 0.649 \\
\bottomrule
\end{tabular}
\end{table}

Three findings emerge. \textbf{(1) Encoder adaptation gives a modest, real
lift.} The best single-model configurations reach a per-seed mean AUC of
$\approx\!0.915$ against the frozen control's $0.892$ --- a $\approx\!+0.02$ AUC
gain that clears the pre-registered adoption band. The lift is smaller and the
field tighter than the ensemble metric suggests: the top LoRA, \texttt{last-}4,
and full configurations ($0.906$--$0.916$) are all within seed noise of one
another, so it does \emph{not} establish a single method as the winner
--- only that lightweight adaptation helps and that the choice among
parameter-efficient methods must be settled on the test split. We nonetheless
note that LoRA supplies the most \emph{stable} top configurations (low per-seed
variance at high mean), whereas the configurations the ensemble metric would
have promoted --- LoRA $r$32/$r$16 at lr1e-4 --- carry $\pm.05$ seed variance and
sit at single-model rank $\approx\!11$, a cautionary example of the
ensemble--single-model mismatch. \textbf{(2) The encoder swap is neutral.}
\texttt{mxbai} matches \texttt{bge} under both frozen ($0.886$ vs.\ $0.892$) and
LoRA ($0.916$ vs.\ $0.915$); it is the nominal rank-1 single model but is within
noise of the \texttt{bge} leaders. \textbf{(3) DAPT collapses.} Both DAPT arms
fall to $0.65$--$0.71$ AUC --- far \emph{below} the no-adaptation baseline ---
and predict the majority class (per-seed F1 $=0$ on several seeds). The
masked-LM objective, applied to a 107-interview corpus with no
embedding-preserving term, degrades the very sentence-embedding geometry the head
pools over: a catastrophic-forgetting failure, not an under-training one.
Lengthening the masked-LM schedule would be expected to worsen it; reviving DAPT
would need an embedding-preserving objective (e.g.\ TSDAE or contrastive
continued pre-training), which we do not pursue.

The pre-registered rule promotes the four configurations within the seed-noise
band of the best (capped at four) --- \texttt{mxbai} LoRA $r$16,
\texttt{bge} LoRA $r$16+FFN, \texttt{bge} LoRA $r$8/lr2e-4, and \texttt{last-}4 ---
plus the frozen control to Stage 2, and carries the rank-1 \texttt{mxbai} LoRA
$r$16 forward as the single winner for the Stage-3 cross-validation.

\subsection{The leakage-free OOF probe}
\label{sec:finetune-iter1-oof}

Before the once-only test evaluation, each finalist passes through the Stage-2
out-of-fold (OOF) probe: a nested $K$-fold over the train+dev pool that gives
every one of the $N\!=\!142$ subjects a leakage-free probability from a model
that never trained on it. The probe exists to fix the decision threshold, but it
is \emph{also} a leakage-free, fourfold-larger-$N$ read on the finalist ranking
than the $n\!=\!35$ development split, and it is informative on its own.
\Cref{tab:ft-oof} reports it.

\begin{table}[t]
\caption{Stage-2 OOF probe (leakage-free, train+dev pool $N\!=\!142$, 42
positive). \textbf{OOF AUC} with a subject-resampling bootstrap 95\% CI; the
development and OOF ranks are repeated for contrast (development per-seed AUC in
\Cref{tab:ft-grid}).}
\label{tab:ft-oof}
\centering\small
\setlength{\tabcolsep}{4pt}
\begin{tabular}{@{}lcccc@{}}
\toprule
\textbf{Model} & \textbf{OOF AUC} & \textbf{95\% CI} & \textbf{dev rk} & \textbf{OOF rk} \\
\midrule
\textbf{frozen (control)}      & \textbf{0.790} & [.709,.866] & 5 & \textbf{1} \\
last-4                         & 0.789 & [.706,.862] & 4 & 2 \\
LoRA $r$16+FFN                 & 0.775 & [.687,.851] & 2 & 3 \\
mxbai LoRA $r$16               & 0.768 & [.685,.843] & \textbf{1} & 4 \\
LoRA $r$8                      & 0.764 & [.676,.846] & 3 & 5 \\
\bottomrule
\end{tabular}
\end{table}

Three points follow, all consistent with the development-grid reading.
\textbf{(1) The grid ranking does not survive the leakage-free probe.} The
development rank-1 configuration (\texttt{mxbai} LoRA $r$16) falls to OOF rank~4,
while \texttt{last-}4 --- development rank~4 --- rises to OOF rank~2; yet every
pairwise OOF AUC gap lies within bootstrap noise, so the finalists are
\emph{statistically indistinguishable} at $N\!=\!142$. This is the same
conclusion the per-seed grid reached --- no single adaptation method is crowned
--- now corroborated on a leakage-free, larger sample, and it is
exactly why selection was pre-registered and the test split is spent once.
\textbf{(2) The development AUC is optimistic.} The leakage-free OOF AUC
($\approx\!0.76$--$0.79$) sits about $0.13$ below the development figure
($\approx\!0.91$); test expectations should be calibrated to the OOF band.
\textbf{(3) Fine-tuning yields no leakage-free OOF lift over the frozen
control.} Run through the identical OOF probe, the frozen encoder scores the
\emph{highest} pooled OOF AUC of the set ($0.790$, $95\%$ CI $[.709,.866]$),
edging the best adaptation arm (\texttt{last-}4, $0.789$). Treating the 15 paired
per-run AUCs ($5$ fold $\times\,3$ seed, aligned by fold and seed) with a
Wilcoxon signed-rank test, no fine-tuning finalist differs from frozen: the
mean per-run gaps are within $\pm0.01$ AUC and the raw $p$-values are
$0.98$/$0.39$/$0.60$/$0.69$ (\texttt{last-}4/LoRA $r$16+FFN/\texttt{mxbai}/LoRA
$r$8), all $\approx\!1$ after Holm correction. Every model separates the pooled
pool far above chance (Mann-Whitney $p<10^{-6}$), but adaptation moves the
\emph{ranking} no further than frozen does. The grid's apparent $\approx\!+0.02$
development-AUC lift over the frozen control is therefore a small-$n$ artefact of
the $n\!=\!35$ dev split; it does not reproduce on the four-fold-larger,
leakage-free OOF read. This is exactly the kind of optimism the staged protocol
was built to catch before spending the test split.

\subsection{Official-test and cross-validation results}
\label{sec:finetune-iter1-test}

The leakage-free OOF probe predicted that the development lift would not survive.
The two once-only readings --- the official AVEC2017 test split and the
cross-validation pool --- confirm it.

\paragraph{Official test split.} Each finalist is evaluated once on the official
test split ($n\!=\!47$, 14 positive) as a 30-seed probability ensemble, with the
decision threshold fixed by the leakage-free OOF prevalence rule of
\Cref{sec:finetune-iter1-oof}; \Cref{tab:ft-test} reports the result against the
frozen single-model bar (AUC $0.864$, macro-F1 $0.774$; lever~3 of
\Cref{tab:levers}). \emph{No fine-tuning finalist beats the frozen bar.} The best
ROC-AUC --- \texttt{last-}4 at $0.857$ --- still sits below frozen, and the
remaining three span $0.816$--$0.846$; every macro-F1 ($0.743$--$0.763$) is below
the frozen $0.774$. The comparison is in fact conservative against frozen: these
are 30-seed \emph{ensemble} predictions, yet they fail to clear a \emph{single}
frozen model. Each AUC bootstrap CI is wide ($\approx\!0.24$ span at $n\!=\!47$)
and contains the frozen bar, so no finalist is \emph{distinguishable} from frozen
either way --- the honest reading is a null, not a regression. The
development-grid ordering again fails to transfer: the dev rank-1 model
(\texttt{mxbai} LoRA $r$16) posts the \emph{lowest} test AUC ($0.816$), while the
dev rank-4 \texttt{last-}4 posts the highest --- the same rank inversion the OOF
probe showed, now reproduced on the test split.

\begin{table}[t]
\caption{Official AVEC2017 test split ($n\!=\!47$, 14 positive), 30-seed
probability ensemble at the leakage-free OOF prevalence-matched threshold.
\textbf{AUC} carries a subject-resampling bootstrap 95\% CI. Top row is the frozen
single-model bar (AUC $0.864$, macro-F1 $0.774$). Rows ordered by test AUC; the
development rank is repeated for contrast.}
\label{tab:ft-test}
\centering\small
\setlength{\tabcolsep}{4pt}
\begin{tabular}{@{}lcccc@{}}
\toprule
\textbf{Model} & \textbf{AUC} & \textbf{95\% CI} & \textbf{mac-F1} & \textbf{dev rk} \\
\midrule
frozen (bar)                   & \textbf{0.864} & --- & \textbf{0.774} & --- \\
last-4, lr1e-5                 & 0.857 & [.722,.964] & 0.743 & 4 \\
LoRA $r$16+FFN, lr1e-4         & 0.846 & [.706,.959] & 0.743 & 2 \\
LoRA $r$8, lr2e-4              & 0.844 & [.708,.953] & 0.743 & 3 \\
mxbai LoRA $r$16, lr1e-4       & 0.816 & [.676,.933] & 0.763 & 1 \\
\bottomrule
\end{tabular}
\end{table}

\paragraph{Cross-validation.} Stage~3 re-evaluates the single winner
(\texttt{mxbai} LoRA $r$16, carried forward as the nominal dev rank-1) against the
frozen control under repeated $K$-fold cross-validation on the train+dev pool, all
adaptation performed inside each fold (\Cref{tab:ft-cv}). Over the 50 paired
per-run AUCs (2 repeats $\times\,5$ fold $\times\,5$ seed, aligned by repeat, fold,
and seed), the winner and frozen are indistinguishable on AUC ($0.777$ vs.\ $0.778$;
Wilcoxon $p=0.79$) and frozen is nominally ahead on macro-F1 ($0.679$ vs.\ $0.657$;
$p=0.12$, n.s.). Monte-Carlo cross-validation, run for both arms,
\emph{flips the sign} of the nominal gap --- the winner
edges frozen on MC (per-run AUC $0.824$ vs.\ $0.803$, $\Delta\!=\!+0.021$,
$p=0.35$; macro-F1 $0.694$ vs.\ $0.667$, $\Delta\!=\!+0.026$, $p=0.47$) having
trailed it on $K$-fold --- but neither gap is significant, and a lift that
reverses direction between two cross-validation schemes is noise, not signal.

\begin{table}[t]
\caption{Stage-3 cross-validation (train+dev pool), winner (\texttt{mxbai} LoRA
$r$16) vs.\ the in-harness frozen control, adaptation inside each fold. Repeated
$K$-fold is $2$ repeat $\times\,5$ fold $\times\,5$ seed $=50$ per-run
evaluations; Monte-Carlo is $25$ per-run. Per-run mean$\pm$std and the
fold-ensemble.}
\label{tab:ft-cv}
\centering\small
\setlength{\tabcolsep}{4pt}
\begin{tabular}{@{}llcccc@{}}
\toprule
 & & \multicolumn{2}{c}{\textbf{per-run AUC}} & \multicolumn{2}{c}{\textbf{per-run mac-F1}} \\
\cmidrule(lr){3-4}\cmidrule(lr){5-6}
\textbf{Scheme} & \textbf{Model} & \textbf{run} & \textbf{ens.} & \textbf{run} & \textbf{ens.} \\
\midrule
$K$-fold & frozen        & \textbf{0.778}{\footnotesize$\pm$.11} & 0.80 & \textbf{0.679}{\footnotesize$\pm$.14} & 0.72 \\
$K$-fold & winner        & 0.777{\footnotesize$\pm$.11} & 0.80 & 0.657{\footnotesize$\pm$.11} & 0.67 \\
\midrule
MC       & frozen        & 0.803{\footnotesize$\pm$.09} & 0.843 & 0.667{\footnotesize$\pm$.12} & 0.732 \\
MC       & winner        & \textbf{0.824}{\footnotesize$\pm$.09} & \textbf{0.847} & \textbf{0.694}{\footnotesize$\pm$.13} & \textbf{0.749} \\
\bottomrule
\end{tabular}
\end{table}

\paragraph{Verdict.} Three independent leakage-free readings --- the
$N\!=\!142$ OOF probe, the once-only $n\!=\!47$ test split, and the repeated
cross-validation pool --- agree: encoder fine-tuning yields \emph{no} lift over the
frozen baseline for this task at this scale. The $\approx\!+0.02$ AUC advantage on
the $n\!=\!35$ development grid was a small-sample artefact; it vanishes the moment
it is measured on a larger, leakage-free sample, and the dev ranking among
adaptation methods never transfers (the dev rank-1 \texttt{mxbai} LoRA $r$16 is the
worst finalist at test). The pre-registered, staged protocol did exactly what it
was designed to do: it exposed the development optimism \emph{before} the test
split was spent, and prevents us from over-claiming a parameter-efficient
fine-tuning win that the data does not support. We therefore retain the frozen
encoder as the deployed system; the remaining headroom (\Cref{tab:levers}) is not
recoverable by fine-tuning the sentence encoder on a 107-interview corpus.

% placed AFTER \section{Internal Comparison: Cross-Validation}
% (\label{sec:internal}) and BEFORE \section{Discussion and Limitations}.
% Iteration 1 is complete on the development grid; the official-test and
% cross-validation blocks marked %TODO are filled once cluster jobs 3/4 land.
% =====================================================================
\section{Does Fine-Tuning the Encoder Help?}
\label{sec:finetune}

Every result so far keeps the sentence encoder \emph{frozen}: the
single-lever study (\Cref{tab:levers}) showed the remaining headroom is in the
\emph{ranking}, not the decision rule, therefore to lift the AUC ceiling a possibility is
to fine-tune the encoder.
Fine-tuning an $n\!=\!107$ corpus is the regime where over-fitting is most
likely, so the study is run as a pre-registered, staged search that selects on
development data and touches the official test split only once, for a small set
of finalists.

\subsection{Fine-tuning menu and staged protocol}
\label{sec:finetune-protocol}
We compare five adaptation regimes against the frozen baseline, ordered by
trainable-parameter budget: (i) \textbf{frozen} (control, embeddings only);
(ii) \textbf{BitFit} (bias terms only); (iii) \textbf{LoRA} low-rank adapters on
the attention (optionally FFN) projections, rank $r\!\in\!\{8,16,32\}$;
(iv) \textbf{last-$k$} unfreezing of the top $k\!\in\!\{2,4\}$ transformer
blocks with layer-wise LR decay; and (v) \textbf{full} fine-tuning. We also test
\textbf{domain-adaptive pre-training} (DAPT): a masked-LM pass over the training
interviews before the supervised stage. The encoder is \texttt{bge-large-en-v1.5}
throughout; one arm swaps in \texttt{mxbai-embed-large-v1} to check
encoder sensitivity. The temporal GRU pooling head
(\Cref{sec:method}) is unchanged.

Selection is staged to ring-fence the test split:
\begin{itemize}
  \item \textbf{Stage 1 --- development grid.} 20 configurations, 5 seeds each,
        \emph{official protocol but no test evaluation}. Configurations are ranked on the 5-seed
        \emph{per-seed mean} development ROC-AUC --- the same quantity reported
        at test time.
  \item \textbf{Stage 2 --- finalists.} Every configuration whose per-seed mean
        development AUC is within one seed-noise band ($\pm0.015$) of the best,
        capped at four, plus the frozen control, advances. Each gets a
        leakage-free decision threshold from a nested out-of-fold (OOF) probe and
        is then evaluated \emph{once} on the official test split over 30 seeds.
  \item \textbf{Stage 3 --- cross-validation.} The single winner and the frozen
        control are re-evaluated under $K$-fold (with a repeat) and
        Monte-Carlo cross-validation on the training pool, mirroring
        \Cref{sec:internal}.
\end{itemize}
All encoder adaptation happens \emph{inside} each fold; no test or
held-out interview is ever seen during fine-tuning or thresholding.

\subsection{Development-grid selection}
\label{sec:finetune-iter1}

A practical caveat governs how the grid is read. The development split is
$n\!=\!35$ interviews (12 positive); at this size the \emph{thresholded} metrics
are quantised --- many unrelated configurations land on the identical confusion
matrix --- so they carry little ranking signal. We therefore rank on the
per-seed mean ROC-AUC, the only continuous discriminator, and treat AUC gaps
below $\approx\!0.01$ as noise. \Cref{tab:ft-grid} reports the representative
configurations.

These dev figures are produced by the fine-tuning harness --- which trains the
MIL head with a shorter, more strongly regularised schedule than the
design-ablation pipeline of \Cref{tab:abldesign-rep} (and embeds under mixed
precision) --- so they are internally consistent but \emph{not} numerically
comparable to that table's frozen dev AUC. All fine-tuning effects are therefore
measured against the in-harness frozen control, not against the design-ablation
numbers; the headline frozen reference remains the once-only official-test bar.

\begin{table}[t]
\caption{Fine-tuning grid, development split, ranked by the
selection metric. \textbf{AUC} is the per-seed mean$\pm$std over 5 seeds (the
selection metric); \textbf{ens.\ AUC} is the 5-seed probability ensemble. Adopted
Stage-2 finalists in \textbf{bold}.}
\label{tab:ft-grid}
\centering
\small
\setlength{\tabcolsep}{4pt}
\begin{tabular}{@{}lcc@{}}
\toprule
\textbf{Configuration} & \textbf{AUC (per-seed)} & \textbf{ens.\ AUC} \\
\midrule
\textbf{mxbai, LoRA $r$16, lr1e-4}  & \textbf{0.916}{\footnotesize$\pm$.005} & 0.917 \\
\textbf{LoRA $r$16+FFN, lr1e-4}     & \textbf{0.915}{\footnotesize$\pm$.008} & 0.920 \\
\textbf{LoRA $r$8, lr2e-4}          & \textbf{0.913}{\footnotesize$\pm$.010} & 0.924 \\
\textbf{last-4, lr1e-5}             & \textbf{0.908}{\footnotesize$\pm$.008} & 0.899 \\
LoRA $r$16, lr2e-4                  & 0.907{\footnotesize$\pm$.022} & 0.920 \\
full, lr1e-5                        & 0.906{\footnotesize$\pm$.006} & 0.899 \\
last-2, lr2e-5                      & 0.899{\footnotesize$\pm$.005} & 0.899 \\
BitFit, lr1e-4                      & 0.899{\footnotesize$\pm$.010} & 0.902 \\
LoRA $r$32, lr1e-4                  & 0.893{\footnotesize$\pm$.053} & 0.924 \\
frozen (control)                   & 0.892{\footnotesize$\pm$.003} & 0.884 \\
mxbai, frozen                      & 0.886{\footnotesize$\pm$.003} & 0.884 \\
LoRA $r$16, lr1e-4                  & 0.888{\footnotesize$\pm$.052} & 0.920 \\
\midrule
DAPT, frozen                       & 0.709{\footnotesize$\pm$.044} & 0.707 \\
DAPT, LoRA $r$16                    & 0.708{\footnotesize$\pm$.078} & 0.649 \\
\bottomrule
\end{tabular}
\end{table}

Three findings emerge. \textbf{(1) Encoder adaptation gives a modest, real
lift.} The best single-model configurations reach a per-seed mean AUC of
$\approx\!0.915$ against the frozen control's $0.892$ --- a $\approx\!+0.02$ AUC
gain that clears the pre-registered adoption band. The lift is smaller and the
field tighter than the ensemble metric suggests: the top LoRA, \texttt{last-}4,
and full configurations ($0.906$--$0.916$) are all within seed noise of one
another, so it does \emph{not} establish a single method as the winner
--- only that lightweight adaptation helps and that the choice among
parameter-efficient methods must be settled on the test split. We nonetheless
note that LoRA supplies the most \emph{stable} top configurations (low per-seed
variance at high mean), whereas the configurations the ensemble metric would
have promoted --- LoRA $r$32/$r$16 at lr1e-4 --- carry $\pm.05$ seed variance and
sit at single-model rank $\approx\!11$, a cautionary example of the
ensemble--single-model mismatch. \textbf{(2) The encoder swap is neutral.}
\texttt{mxbai} matches \texttt{bge} under both frozen ($0.886$ vs.\ $0.892$) and
LoRA ($0.916$ vs.\ $0.915$); it is the nominal rank-1 single model but is within
noise of the \texttt{bge} leaders. \textbf{(3) DAPT collapses.} Both DAPT arms
fall to $0.65$--$0.71$ AUC --- far \emph{below} the no-adaptation baseline ---
and predict the majority class (per-seed F1 $=0$ on several seeds). The
masked-LM objective, applied to a 107-interview corpus with no
embedding-preserving term, degrades the very sentence-embedding geometry the head
pools over: a catastrophic-forgetting failure, not an under-training one.
Lengthening the masked-LM schedule would be expected to worsen it; reviving DAPT
would need an embedding-preserving objective (e.g.\ TSDAE or contrastive
continued pre-training), which we do not pursue.

The pre-registered rule promotes the four configurations within the seed-noise
band of the best (capped at four) --- \texttt{mxbai} LoRA $r$16,
\texttt{bge} LoRA $r$16+FFN, \texttt{bge} LoRA $r$8/lr2e-4, and \texttt{last-}4 ---
plus the frozen control to Stage 2, and carries the rank-1 \texttt{mxbai} LoRA
$r$16 forward as the single winner for the Stage-3 cross-validation.

\subsection{The leakage-free OOF probe}
\label{sec:finetune-iter1-oof}

Before the once-only test evaluation, each finalist passes through the Stage-2
out-of-fold (OOF) probe: a nested $K$-fold over the train+dev pool that gives
every one of the $N\!=\!142$ subjects a leakage-free probability from a model
that never trained on it. The probe exists to fix the decision threshold, but it
is \emph{also} a leakage-free, fourfold-larger-$N$ read on the finalist ranking
than the $n\!=\!35$ development split, and it is informative on its own.
\Cref{tab:ft-oof} reports it.

\begin{table}[t]
\caption{Stage-2 OOF probe (leakage-free, train+dev pool $N\!=\!142$, 42
positive). \textbf{OOF AUC} with a subject-resampling bootstrap 95\% CI; the
development and OOF ranks are repeated for contrast (development per-seed AUC in
\Cref{tab:ft-grid}).}
\label{tab:ft-oof}
\centering\small
\setlength{\tabcolsep}{4pt}
\begin{tabular}{@{}lcccc@{}}
\toprule
\textbf{Model} & \textbf{OOF AUC} & \textbf{95\% CI} & \textbf{dev rk} & \textbf{OOF rk} \\
\midrule
\textbf{frozen (control)}      & \textbf{0.790} & [.709,.866] & 5 & \textbf{1} \\
last-4                         & 0.789 & [.706,.862] & 4 & 2 \\
LoRA $r$16+FFN                 & 0.775 & [.687,.851] & 2 & 3 \\
mxbai LoRA $r$16               & 0.768 & [.685,.843] & \textbf{1} & 4 \\
LoRA $r$8                      & 0.764 & [.676,.846] & 3 & 5 \\
\bottomrule
\end{tabular}
\end{table}

Three points follow, all consistent with the development-grid reading.
\textbf{(1) The grid ranking does not survive the leakage-free probe.} The
development rank-1 configuration (\texttt{mxbai} LoRA $r$16) falls to OOF rank~4,
while \texttt{last-}4 --- development rank~4 --- rises to OOF rank~2; yet every
pairwise OOF AUC gap lies within bootstrap noise, so the finalists are
\emph{statistically indistinguishable} at $N\!=\!142$. This is the same
conclusion the per-seed grid reached --- no single adaptation method is crowned
--- now corroborated on a leakage-free, larger sample, and it is
exactly why selection was pre-registered and the test split is spent once.
\textbf{(2) The development AUC is optimistic.} The leakage-free OOF AUC
($\approx\!0.76$--$0.79$) sits about $0.13$ below the development figure
($\approx\!0.91$); test expectations should be calibrated to the OOF band.
\textbf{(3) Fine-tuning yields no leakage-free OOF lift over the frozen
control.} Run through the identical OOF probe, the frozen encoder scores the
\emph{highest} pooled OOF AUC of the set ($0.790$, $95\%$ CI $[.709,.866]$),
edging the best adaptation arm (\texttt{last-}4, $0.789$). Treating the 15 paired
per-run AUCs ($5$ fold $\times\,3$ seed, aligned by fold and seed) with a
Wilcoxon signed-rank test, no fine-tuning finalist differs from frozen: the
mean per-run gaps are within $\pm0.01$ AUC and the raw $p$-values are
$0.98$/$0.39$/$0.60$/$0.69$ (\texttt{last-}4/LoRA $r$16+FFN/\texttt{mxbai}/LoRA
$r$8), all $\approx\!1$ after Holm correction. Every model separates the pooled
pool far above chance (Mann-Whitney $p<10^{-6}$), but adaptation moves the
\emph{ranking} no further than frozen does. The grid's apparent $\approx\!+0.02$
development-AUC lift over the frozen control is therefore a small-$n$ artefact of
the $n\!=\!35$ dev split; it does not reproduce on the four-fold-larger,
leakage-free OOF read. This is exactly the kind of optimism the staged protocol
was built to catch before spending the test split.

\subsection{Official-test and cross-validation results}
\label{sec:finetune-iter1-test}

The leakage-free OOF probe predicted that the development lift would not survive.
The two once-only readings --- the official AVEC2017 test split and the
cross-validation pool --- confirm it.

\paragraph{Official test split.} Each finalist is evaluated once on the official
test split ($n\!=\!47$, 14 positive) as a 30-seed probability ensemble, with the
decision threshold fixed by the leakage-free OOF prevalence rule of
\Cref{sec:finetune-iter1-oof}; \Cref{tab:ft-test} reports the result against the
frozen single-model bar (AUC $0.864$, macro-F1 $0.774$; lever~3 of
\Cref{tab:levers}). \emph{No fine-tuning finalist beats the frozen bar.} The best
ROC-AUC --- \texttt{last-}4 at $0.857$ --- still sits below frozen, and the
remaining three span $0.816$--$0.846$; every macro-F1 ($0.743$--$0.763$) is below
the frozen $0.774$. The comparison is in fact conservative against frozen: these
are 30-seed \emph{ensemble} predictions, yet they fail to clear a \emph{single}
frozen model. Each AUC bootstrap CI is wide ($\approx\!0.24$ span at $n\!=\!47$)
and contains the frozen bar, so no finalist is \emph{distinguishable} from frozen
either way --- the honest reading is a null, not a regression. The
development-grid ordering again fails to transfer: the dev rank-1 model
(\texttt{mxbai} LoRA $r$16) posts the \emph{lowest} test AUC ($0.816$), while the
dev rank-4 \texttt{last-}4 posts the highest --- the same rank inversion the OOF
probe showed, now reproduced on the test split.

\begin{table}[t]
\caption{Official AVEC2017 test split ($n\!=\!47$, 14 positive), 30-seed
probability ensemble at the leakage-free OOF prevalence-matched threshold.
\textbf{AUC} carries a subject-resampling bootstrap 95\% CI. Top row is the frozen
single-model bar (AUC $0.864$, macro-F1 $0.774$). Rows ordered by test AUC; the
development rank is repeated for contrast.}
\label{tab:ft-test}
\centering\small
\setlength{\tabcolsep}{4pt}
\begin{tabular}{@{}lcccc@{}}
\toprule
\textbf{Model} & \textbf{AUC} & \textbf{95\% CI} & \textbf{mac-F1} & \textbf{dev rk} \\
\midrule
frozen (bar)                   & \textbf{0.864} & --- & \textbf{0.774} & --- \\
last-4, lr1e-5                 & 0.857 & [.722,.964] & 0.743 & 4 \\
LoRA $r$16+FFN, lr1e-4         & 0.846 & [.706,.959] & 0.743 & 2 \\
LoRA $r$8, lr2e-4              & 0.844 & [.708,.953] & 0.743 & 3 \\
mxbai LoRA $r$16, lr1e-4       & 0.816 & [.676,.933] & 0.763 & 1 \\
\bottomrule
\end{tabular}
\end{table}

\paragraph{Cross-validation.} Stage~3 re-evaluates the single winner
(\texttt{mxbai} LoRA $r$16, carried forward as the nominal dev rank-1) against the
frozen control under repeated $K$-fold cross-validation on the train+dev pool, all
adaptation performed inside each fold (\Cref{tab:ft-cv}). Over the 50 paired
per-run AUCs (2 repeats $\times\,5$ fold $\times\,5$ seed, aligned by repeat, fold,
and seed), the winner and frozen are indistinguishable on AUC ($0.777$ vs.\ $0.778$;
Wilcoxon $p=0.79$) and frozen is nominally ahead on macro-F1 ($0.679$ vs.\ $0.657$;
$p=0.12$, n.s.). Monte-Carlo cross-validation, run for both arms,
\emph{flips the sign} of the nominal gap --- the winner
edges frozen on MC (per-run AUC $0.824$ vs.\ $0.803$, $\Delta\!=\!+0.021$,
$p=0.35$; macro-F1 $0.694$ vs.\ $0.667$, $\Delta\!=\!+0.026$, $p=0.47$) having
trailed it on $K$-fold --- but neither gap is significant, and a lift that
reverses direction between two cross-validation schemes is noise, not signal.

\begin{table}[t]
\caption{Stage-3 cross-validation (train+dev pool), winner (\texttt{mxbai} LoRA
$r$16) vs.\ the in-harness frozen control, adaptation inside each fold. Repeated
$K$-fold is $2$ repeat $\times\,5$ fold $\times\,5$ seed $=50$ per-run
evaluations; Monte-Carlo is $25$ per-run. Per-run mean$\pm$std and the
fold-ensemble.}
\label{tab:ft-cv}
\centering\small
\setlength{\tabcolsep}{4pt}
\begin{tabular}{@{}llcccc@{}}
\toprule
 & & \multicolumn{2}{c}{\textbf{per-run AUC}} & \multicolumn{2}{c}{\textbf{per-run mac-F1}} \\
\cmidrule(lr){3-4}\cmidrule(lr){5-6}
\textbf{Scheme} & \textbf{Model} & \textbf{run} & \textbf{ens.} & \textbf{run} & \textbf{ens.} \\
\midrule
$K$-fold & frozen        & \textbf{0.778}{\footnotesize$\pm$.11} & 0.80 & \textbf{0.679}{\footnotesize$\pm$.14} & 0.72 \\
$K$-fold & winner        & 0.777{\footnotesize$\pm$.11} & 0.80 & 0.657{\footnotesize$\pm$.11} & 0.67 \\
\midrule
MC       & frozen        & 0.803{\footnotesize$\pm$.09} & 0.843 & 0.667{\footnotesize$\pm$.12} & 0.732 \\
MC       & winner        & \textbf{0.824}{\footnotesize$\pm$.09} & \textbf{0.847} & \textbf{0.694}{\footnotesize$\pm$.13} & \textbf{0.749} \\
\bottomrule
\end{tabular}
\end{table}

\paragraph{Verdict.} Three independent leakage-free readings --- the
$N\!=\!142$ OOF probe, the once-only $n\!=\!47$ test split, and the repeated
cross-validation pool --- agree: encoder fine-tuning yields \emph{no} lift over the
frozen baseline for this task at this scale. The $\approx\!+0.02$ AUC advantage on
the $n\!=\!35$ development grid was a small-sample artefact; it vanishes the moment
it is measured on a larger, leakage-free sample, and the dev ranking among
adaptation methods never transfers (the dev rank-1 \texttt{mxbai} LoRA $r$16 is the
worst finalist at test). The pre-registered, staged protocol did exactly what it
was designed to do: it exposed the development optimism \emph{before} the test
split was spent, and prevents us from over-claiming a parameter-efficient
fine-tuning win that the data does not support. We therefore retain the frozen
encoder as the deployed system; the remaining headroom (\Cref{tab:levers}) is not
recoverable by fine-tuning the sentence encoder on a 107-interview corpus.
